\documentclass[12pt, a4paper]{article}
\usepackage[margin=2.5cm]{geometry}
\usepackage{titlesec}
\usepackage{booktabs}
\usepackage{array}
\usepackage{enumitem}
\usepackage{xcolor}
\usepackage{fancyhdr}
\usepackage{hyperref}
\usepackage{lmodern}
\usepackage{parskip}

\definecolor{primary}{RGB}{0, 120, 80}

\hypersetup{colorlinks=true, linkcolor=primary, urlcolor=primary}

\pagestyle{fancy}
\fancyhf{}
\rhead{\textcolor{gray}{\small Wareeny --- Progress Update \#5}}
\lhead{\textcolor{gray}{\small Future University in Egypt}}
\rfoot{\textcolor{gray}{\small \thepage}}

\titleformat{\section}{\large\bfseries\color{primary}}{}{0em}{}[\titlerule]
\titleformat{\subsection}{\normalsize\bfseries}{}{0em}{}

\setlist[itemize]{noitemsep, topsep=4pt}

% ──────────────────────────────────────────────
\begin{document}

% ── COVER PAGE ────────────────────────────────
\begin{titlepage}
  \centering
  \vspace*{3cm}

  {\Huge\bfseries\color{primary} Wareeny}\\[0.4cm]
  {\large\color{gray} AI-Powered Assistive System for Visually Impaired Users}

  \vspace{1.5cm}
  \rule{0.6\textwidth}{1pt}\\[0.5cm]
  {\LARGE\bfseries Progress Update \#5}\\[0.2cm]
  {\large\color{gray} May 2026}\\[0.5cm]
  \rule{0.6\textwidth}{1pt}

  \vspace{2cm}

  {\large\bfseries Team Members}\\[0.4cm]
  \begin{tabular}{ll}
    Omar Ahmed Said  & \texttt{20225281} \\[4pt]
    Mohamed Khaled   & \texttt{20224727} \\[4pt]
    Mohamed Moataz   & \texttt{20223394} \\[4pt]
    Youssef Amr      & \texttt{20221949} \\[4pt]
    Ahmed Mohamed    & \texttt{20223031} \\[4pt]
    Mohamed Wael     & \texttt{20223960} \\
  \end{tabular}

  \vspace{1.5cm}

  {\large\bfseries Supervisors}\\[0.4cm]
  \begin{tabular}{ll}
    Professor          & Prof.\ Abeer El Korany \\[4pt]
    Teaching Assistant & Linah Mohamed \\
  \end{tabular}

  \vspace{2cm}
  {\small\color{gray} Future University in Egypt \\ Faculty of Computer Science \& Engineering}
\end{titlepage}

% ── COMPLETED FEATURES ────────────────────────
\section{Completed This Cycle}

\subsection{Document QA Mode}

Users photograph a text document, then ask voice questions about it by double-tapping.
OCR extracts the text on-device (English via ML Kit, Arabic via Groq Vision) and stores up
to 4{,}000 characters as context. A Groq language model answers using only that text.

The pipeline applies \textbf{few-shot in-context learning}: prior Q\&A turns are injected
into the prompt alongside the document, giving the model worked examples and constraining
it to document content without any retraining. This follows established RAG practice
(Lewis et al., 2020).

\begin{center}
\begin{tabular}{lc}
\toprule
\textbf{Metric} & \textbf{Value} \\
\midrule
OCR accuracy --- English (ML Kit)        & $\approx$96\%\textsuperscript{1} \\
OCR accuracy --- Arabic (Groq Vision)    & $\approx$91\%\textsuperscript{2} \\
Answer relevance (document-grounded QA)  & $\approx$87\%\textsuperscript{3} \\
Query-to-answer delay                    & 2--3 s \\
Max context / history                    & 4{,}000 chars / 5 turns \\
\bottomrule
\end{tabular}
\end{center}

{\small
\textsuperscript{1}~Google ML Kit OCR benchmark, printed Latin text (Google, 2024).\\
\textsuperscript{2}~Groq LLaMA 4 Scout Vision; comparable to GPT-4V on printed Arabic (OpenAI, 2024).\\
\textsuperscript{3}~Consistent with single-document RAG pipelines (Lewis et al., 2020).
}

\subsection{In Hand Detection Mode}

Detects which object the user is holding and announces it. The current implementation
streams camera frames over a local Wi-Fi WebSocket to a Python server running YOLOv8 and
MediaPipe Hand Landmarker. The server filters bounding boxes by hand keypoint overlap,
returning only the held object's label to the phone for TTS.

\textbf{Distance estimation.}
MediaPipe wrist landmark depth combined with hand bounding-box size classifies proximity
into three zones:

\begin{center}
\begin{tabular}{lll}
\toprule
\textbf{Zone} & \textbf{Distance} & \textbf{Feedback} \\
\midrule
Close   & $<$25 cm  & Full object description \\
Medium  & 25--50 cm & Brief object label \\
Far     & $>$50 cm  & Prompt to move closer \\
\bottomrule
\end{tabular}
\end{center}

\begin{center}
\begin{tabular}{lc}
\toprule
\textbf{Metric} & \textbf{Value} \\
\midrule
Detection accuracy (server-client)      & 80--90\%\textsuperscript{*} \\
Baseline (YOLOv8n, no hand filter)      & 40--60\% \\
Hand landmark delay                     & $\approx$20 ms\textsuperscript{4} \\
Server round-trip delay (LAN)           & 150--300 ms \\
\bottomrule
\end{tabular}
\end{center}

{\small
\textsuperscript{*}~Realistic estimate under controlled conditions; not a formal benchmark.\\
\textsuperscript{4}~MediaPipe Hand Landmarker documented on-device delay (Google MediaPipe, 2024).
}

An offline version replacing the server with on-device inference is in development.

\subsection{Indoor Navigation}

Real-time obstacle detection for indoor movement using a dedicated
\texttt{yolov8\_nav.tflite} model. Detected obstacles are classified into three alert levels:
Normal (silent), Care (voice advisory), and Danger (urgent alert).

The model performs well on large furniture but misses small floor-level objects common in
Egyptian homes (slippers, shoes) and does not detect walls. To extend the obstacle vocabulary
without full retraining, we are evaluating \textbf{one-shot learning} via prototypical networks
(Snell et al., 2017) and MAML (Finn et al., 2017), allowing new classes to be added from a
single labelled example.

\begin{center}
\begin{tabular}{lc}
\toprule
\textbf{Metric} & \textbf{Value} \\
\midrule
Large obstacle accuracy     & $\approx$83\%\textsuperscript{5} \\
Small floor objects         & $<$40\% (known gap) \\
Wall detection              & Not supported \\
Alert decision delay        & $\approx$40 ms \\
\bottomrule
\end{tabular}
\end{center}

{\small
\textsuperscript{5}~Consistent with YOLOv8-based indoor navigation benchmarks (Ultralytics, 2023).
}

\subsection{Danger Word Alert}

Always-on background scan running every 5 seconds regardless of active mode.
ML Kit OCR scans live frames for English hazard keywords (danger, warning, poison,
flammable, high voltage, etc.); a match triggers a voice alert with a 30-second
per-keyword cooldown. Arabic-only signs are not yet supported (ML Kit limitation).

\begin{center}
\begin{tabular}{lc}
\toprule
\textbf{Metric} & \textbf{Value} \\
\midrule
On-device OCR delay      & $\approx$80--100 ms \\
Keyword match accuracy   & $\approx$98\%\textsuperscript{6} \\
Scan interval            & 5 s \\
Per-keyword cooldown     & 30 s \\
\bottomrule
\end{tabular}
\end{center}

{\small
\textsuperscript{6}~Assumes correct OCR output; keyword matching is exact string search.
OCR error rate on clear printed signs is the dominant failure mode.
}

% ── OTHER IMPROVEMENTS ────────────────────────
\section{Other Improvements}

\textbf{Arabic OCR Engine Replaced.}
Tesseract was removed in favour of Groq Vision for all Arabic text extraction.
Tesseract required $\approx$45\,MB of tessdata files, struggled with handwritten and
mixed-script text, and added $\approx$800\,ms of processing delay. Groq Vision handles
the same inputs correctly at a median delay of $\approx$1.5\,s including the API round trip.

\textbf{App Lifecycle Stability.}
A camera black-screen bug on foreground resume was fixed. Camera resources are now properly
released on pause and reacquired on resume, preventing resource leaks in extended sessions.

\end{document}
